\documentclass[a4paper,11pt]{article}
\usepackage[utf8]{inputenc}
\usepackage{amsmath, amssymb}
\usepackage{graphicx}
\usepackage{hyperref}
\usepackage{listings}
\usepackage{xcolor}
\usepackage{geometry}
\geometry{margin=2.5cm}

\title{\textbf{PixHomology Benchmark: User and Reproducibility Guide}}
\author{PixHomology Development Team}
\date{\today}

\begin{document}
\maketitle

\tableofcontents
\newpage

\section{Introduction}
PixHomology is a C++/pybind11 library for computing persistent homology on 2D images. It provides an efficient and memory-aware implementation for computing 0-dimensional and 1-dimensional persistence diagrams from grayscale or binary images represented as NumPy 2D arrays. The accompanying \texttt{PHBenchmark} package automates the benchmarking process across multiple persistent homology backends (PixHomology, Cripser, GUDHI, Ripser, TTK) and datasets.

This document explains:
\begin{itemize}
    \item The internal logic and basic usage of PixHomology.
    \item How to compute persistent homology from NumPy arrays.
    \item How to install and use the PHBenchmark repository for reproducing experiments and reports.
\end{itemize}

\section{PixHomology: Core Functionality}

\subsection{Basic Workflow}
PixHomology converts a 2D NumPy array \texttt{img} into a complex and computes persistence diagrams through lower-star filtrations.

The typical steps are:
\begin{enumerate}
    \item Import the library and provide a 2D NumPy array.
    \item Choose the homology dimension (0 or 1).
    \item Compute the persistence diagram.
\end{enumerate}

\subsection{Example Usage}
\begin{lstlisting}[language=Python, caption=Basic persistent homology computation using PixHomology]
import numpy as np
import pixhomology as px

# Example: create a synthetic 2D image
img = np.random.rand(28, 28)

# Compute 0- and 1-dimensional persistence diagrams
ph0 = px.computePH(img, maxdim=0)
ph0, ph1 = px.computePH(img, maxdim=1)

print("H0 diagram:", ph0)
print("H1 diagram:", ph1)
\end{lstlisting}

The output is a list of (birth, death) pairs for each topological feature. The function automatically handles grayscale-to-filtration conversion using lower-star filtration rules.

\subsection{Documentation and API}
For a complete description of all functionalities and the full API of PixHomology, please refer to the official library documentation available at:

https://riccardoc95.github.io/PixHomology/

\section{Repository Overview}
The \texttt{PHBenchmark} project wraps PixHomology and other libraries into a unified benchmarking toolchain.

\subsection{Main Structure}
\begin{verbatim}
PHBenchmark/
 environment.yml               # Conda environment for full reproducibility
 run_all.sbatch                # SLURM array job for distributed benchmarking
 phbenchmark/
    app.py                     # CLI entry point (Typer)
    download.py                # dataset fetching/extraction
    run.py                     # executes benchmarks via /usr/bin/time
    collect.py                 # aggregates CSV -> JSON summaries
    report.py                  # generates plots and LaTeX tables
    tests/                     # wrappers for each method (pixh, gudhi, etc.)
 PixHomology/                  # local source for PixHomology library
\end{verbatim}

\section{Installation}
The repository includes an \texttt{environment.yml} file that builds a complete reproducible environment.

\subsection{Creating the Environment}
\begin{lstlisting}[language=bash]
conda env create -f environment.yml
conda activate phbenchmark
\end{lstlisting}

This installs:
\begin{itemize}
    \item Python 3.11 and Topology ToolKit (TTK)
    \item Persistent homology libraries: Ripser, Cripser, GUDHI
    \item PixHomology (editable install from local source)
    \item PHBenchmark CLI
\end{itemize}

\subsection{Verification}
\begin{lstlisting}[language=bash]
phbenchmark --help
python -c "import pixhomology as px; print('PixHomology OK', px.__version__)"
\end{lstlisting}

\section{Using PHBenchmark}

\subsection{Dataset Download}
\begin{lstlisting}[language=bash]
phbenchmark download --datasets_dir datasets --dataset mnist
\end{lstlisting}

Available datasets: MNIST, CIFAR10, ImageNet, DIV2K, Kather. Images are expanded to \texttt{datasets/<name>\_npy/}.

\subsection{Running Benchmarks}
Each run measures runtime and peak memory usage.
\begin{lstlisting}[language=bash]
phbenchmark run \
  --method pixh \
  --dataset datasets/mnist_npy \
  --maxdim 1 \
  --output results/pixh/mnist/dim1.csv
\end{lstlisting}

\subsection{Collecting and Reporting}
\begin{lstlisting}[language=bash]
phbenchmark collect \
  --results results \
  --output results/benchmarks.json \
  --summary results/benchmarks_summary.json

phbenchmark report \
  --summary results/benchmarks_summary.json \
  --plots-dir results/plots \
  --latex-output results/benchmark_table.tex
\end{lstlisting}

Outputs include detailed plots and a LaTeX table summarizing average runtime and memory per method and dataset.

\section{End-to-End MNIST Example}
\begin{lstlisting}[language=bash]
conda activate phbenchmark
phbenchmark download --datasets_dir datasets --dataset mnist

for m in pixh cripser gudhi ripser ttk; do
  for k in 0 1; do
    out="results/${m}/mnist/dim${k}.csv"
    mkdir -p $(dirname $out)
    phbenchmark run --method $m --dataset datasets/mnist_npy \
        --maxdim $k --output $out
  done
done

phbenchmark collect --results results --output results/benchmarks.json \
    --summary results/benchmarks_summary.json
phbenchmark report --summary results/benchmarks_summary.json \
    --plots-dir results/plots --latex-output results/benchmark_table.tex
\end{lstlisting}

\section{Running on SLURM}
The provided \texttt{run\_all.sbatch} automates all runs as a SLURM job array.
\begin{lstlisting}[language=bash]
sbatch --array=0-<N-1> run_all.sbatch
\end{lstlisting}

Edit the file to configure datasets, methods, and dimensions. Logs are stored under \texttt{logs/}.

\section{Software Installed}
\begin{itemize}
    \item \textbf{PixHomology} - This repository (C++/pybind11)
    \item \textbf{Cripser} - \url{https://github.com/shizuo-kaji/CubicalRipser_3dim}
    \item \textbf{GUDHI} - \url{https://gudhi.inria.fr/}
    \item \textbf{Ripser.py} - \url{https://ripser.scikit-tda.org/en/latest/}
    \item \textbf{Topology ToolKit (TTK)} - \url{https://topology-tool-kit.github.io/}
\end{itemize}

General Python utilities: \texttt{typer}, \texttt{tqdm}, \texttt{pandas}, \texttt{matplotlib}, \texttt{seaborn}, \texttt{jinja2}.

\section{Reproducibility Notes}
\begin{itemize}
    \item Versions pinned in \texttt{environment.yml}.
    \item Runtime and memory measured using \texttt{/usr/bin/time}.
    \item CPU threads controlled via \texttt{OMP\_NUM\_THREADS}.
\end{itemize}

\end{document}

